% ===========================================================================
% Chapter 2 -- Mathematical Foundations
% Rebuilt in the v3 didactic standard. Organised by what machine learning
% actually uses, in dependency order, with a worked numeric example for
% every mechanism.
% ===========================================================================

\chapter{Mathematical Foundations}
\label{ch:mathematical-foundations}

\chapterlead{Which mathematical objects make learning, uncertainty, and
optimisation precise?}{Objects}{Transform}{Estimate}{Optimise}{mathematics}

\begin{learningobjectives}
After working through this chapter you should be able to:
\begin{itemize}
  \item state the shape of every object in an expression before evaluating it,
        and use a shape mismatch as a debugging tool;
  \item compute the $L_1$, $L_2$, and $L_\infty$ norms of a vector and say which
        one a given application needs;
  \item multiply matrices by hand and predict the output shape from the inputs;
  \item find the eigenvalues of a $2\times2$ matrix and check them against the
        trace and determinant;
  \item apply the chain rule to a composed function --- the single result that
        makes backpropagation possible;
  \item run three steps of gradient descent by hand and derive the learning
        rate at which it diverges;
  \item update a belief with Bayes' theorem and explain why a $99\,\%$ accurate
        test can still be wrong half the time;
  \item compute entropy, cross-entropy, and KL divergence, and show that KL is
        not symmetric;
  \item say what a $p$-value is not;
  \item recognise when a result is mathematically correct and numerically
        worthless.
\end{itemize}
\end{learningobjectives}

\begin{plainlanguage}
Machine learning uses a surprisingly small amount of mathematics, but it uses
it constantly. Four ideas carry almost everything: a \emph{vector} is a list of
numbers you can treat as a point in space; a \emph{matrix} is a machine that
turns one vector into another; a \emph{derivative} tells you which way to nudge
a number to make an outcome better; and a \emph{probability} is a way of
writing down that you are not certain. Every method later in this book is a
particular arrangement of those four things.
\end{plainlanguage}

\begin{engineernote}
You do not need to be fast at algebra to be good at this. You need to be
\emph{unwilling to move on} when a symbol is unclear. Nearly every student who
struggles with deep learning is struggling with a piece of notation nobody
explained --- most often a shape, an index, or a subscript that silently
changed meaning. When you meet a formula in this book, ask three questions
before anything else: what shape is each object, what are its units, and what
happens if I set it to zero.
\end{engineernote}

% ===========================================================================
\section{How to Read Mathematics in This Book}
% ===========================================================================

\subsection{Four Objects, and the Shape Rule}

Almost every expression in this book is built from four kinds of object. The
core notation table in the front matter (\cref{tab:frontmatter-core-notation})
gives the symbols; this section gives the habit.

\begin{figure}[htbp]
\centering
\begin{tikzpicture}[
  lbl/.style={font=\sffamily\scriptsize,color=AlmanachInk!75,align=center},
  hd/.style={font=\sffamily\small\bfseries,color=AlmanachTeal},
  cell/.style={draw=AlmanachTeal!55,fill=AlmanachMist,minimum size=5.5mm,
    inner sep=0pt,font=\scriptsize}]

  \node[hd] at (0,1.5) {scalar};
  \node[cell] at (0,0.6) {$7$};
  \node[lbl] at (0,-0.55) {one number\\\texttt{()}};

  \node[hd] at (2.7,1.5) {vector};
  \foreach \i in {0,1,2}{\node[cell] at (2.7,1.0-\i*5.5mm) {};}
  \node[lbl] at (2.7,-0.55) {a point in space\\\texttt{(d,)}};

  \node[hd] at (5.6,1.5) {matrix};
  \foreach \i in {0,1,2}\foreach \j in {0,1}{
    \node[cell] at (5.35+\j*5.5mm,1.0-\i*5.5mm) {};}
  \node[lbl] at (5.6,-0.55) {a transformation\\\texttt{(m, n)}};

  \node[hd] at (8.9,1.5) {tensor};
  \foreach \k in {0,1}\foreach \i in {0,1,2}\foreach \j in {0,1}{
    \node[cell,fill=AlmanachOchre!\the\numexpr14+\k*12\relax]
      at (8.5+\j*5.5mm+\k*1.6mm, 1.0-\i*5.5mm+\k*1.6mm) {};}
  \node[lbl] at (8.9,-0.55) {a batch of them\\\texttt{(B, m, n)}};
\end{tikzpicture}
\caption{The four objects, with the shape notation used throughout the book.
The batch dimension $B$ on the right is the one most often forgotten; it is
also the one that causes the most confusing errors, because a wrong batch
dimension frequently still runs.}
\label{fig:math-four-objects}
\end{figure}

\begin{engineernote}[The shape rule]
Before you evaluate anything, write the shape of every object above it. Two
objects can only be added if their shapes match; they can only be multiplied if
the \emph{inner} dimensions match. Roughly nine out of ten errors in
hand-written derivations and in code are shape errors, and a shape error is
visible without doing any arithmetic at all.
\end{engineernote}

\begin{checkpoint}
You are told $\mathbf{y} = W\mathbf{x} + \mathbf{b}$ where
$\mathbf{x}\in\mathbb{R}^{5}$ and $\mathbf{y}\in\mathbb{R}^{3}$. Without any
further information, state the shape of $W$ and of $\mathbf{b}$. Then say what
goes wrong if you accidentally use $W^{\!\top}$.
\end{checkpoint}

% ===========================================================================
\section{Vectors: Data as Geometry}
% ===========================================================================

\subsection{What a Vector Is}

In broad terms, vectors are things one can add, and linear functions are
functions of vectors that respect that
addition\textsuperscript{\cite[][p.~12]{cherneyLinearAlgebra2013}}. Order
matters: a vector is an ordered list, so swapping two components produces a
different object.
\begin{equation}
  \vec{v} =
  \begin{pNiceMatrix}
    a_1 \\ a_2 \\ \vdots \\ a_n
  \end{pNiceMatrix}
  \neq
  \begin{pNiceMatrix}
    \CodeBefore
    \cellcolor[HTML]{FFF3D6}{1-1,2-1}
    \Body
    a_2 \\ a_1 \\ \vdots \\ a_n
  \end{pNiceMatrix}
  \label{eq:math-vector-order}
\end{equation}
\where{$a_1,\dots,a_n$ are the components and $n$ the dimension; the highlighted
cells show that exchanging two of them yields a different vector.}

\begin{plainlanguage}
A vector is a recipe for getting somewhere. ``Three steps east, four steps
north'' is a two-dimensional vector. Everything that follows --- length,
angle, similarity --- is just asking geometric questions about that walk. When
a language model represents the word \emph{cat} as $576$ numbers, it is
describing a walk in $576$ directions at once. You cannot picture that, but
every formula you learn in two dimensions still applies unchanged.
\end{plainlanguage}

\subsection{Length, and the Three Norms You Will Meet}

The \keyterm{norm} measures the size of a vector. The familiar one is the
Euclidean or $2$-norm\textsuperscript{\cite[][p.~90]{cherneyLinearAlgebra2013}},
but machine learning uses three:
\begin{equation}
  \lVert \vec{x}\rVert_1 = \sum_{i=1}^{n}\lvert x_i\rvert,
  \qquad
  \lVert \vec{x}\rVert_2 = \sqrt{\sum_{i=1}^{n} x_i^2},
  \qquad
  \lVert \vec{x}\rVert_\infty = \max_i \lvert x_i\rvert .
  \label{eq:math-norms}
\end{equation}
\where{$x_i$ is the $i$-th component of $\vec{x}$ and $n$ the number of
components; the subscript on $\lVert\cdot\rVert$ names the norm.}

\begin{workedexample}[All three norms of one vector]
Let $\vec{x} = (3,\,-4,\,12)$.
\[
\begin{aligned}
  \lVert\vec{x}\rVert_1 &= \lvert3\rvert+\lvert-4\rvert+\lvert12\rvert
    = 3+4+12 = 19,\\
  \lVert\vec{x}\rVert_2 &= \sqrt{3^2+(-4)^2+12^2}
    = \sqrt{9+16+144} = \sqrt{169} = 13,\\
  \lVert\vec{x}\rVert_\infty &= \max(3,4,12) = 12 .
\end{aligned}
\]
Three different answers for the same vector, and all three are correct --- they
answer different questions. Note the ordering
$\lVert\vec{x}\rVert_\infty \le \lVert\vec{x}\rVert_2 \le
\lVert\vec{x}\rVert_1$, which holds for every vector and is worth remembering
as a sanity check on your arithmetic.
\end{workedexample}

\begin{figure}[htbp]
\centering
\begin{tikzpicture}[scale=1.15,
  ax/.style={->,draw=AlmanachInk!55,>={Stealth[length=1.6mm]}},
  lbl/.style={font=\sffamily\scriptsize,color=AlmanachInk!75}]
  \foreach \dx/\name/\desc in {0/{$L_1$}/{diamond}, 3.4/{$L_2$}/{circle},
      6.8/{$L_\infty$}/{square}}{
    \draw[ax] (\dx-1.25,0) -- (\dx+1.25,0);
    \draw[ax] (\dx,-1.25) -- (\dx,1.25);
    \node[lbl] at (\dx,-1.65) {\name\ unit ball (a \desc)};}
  \draw[thick,AlmanachTeal,fill=AlmanachTeal!12]
    (-1,0) -- (0,1) -- (1,0) -- (0,-1) -- cycle;
  \draw[thick,AlmanachTeal,fill=AlmanachTeal!12] (3.4,0) circle (1);
  \draw[thick,AlmanachTeal,fill=AlmanachTeal!12]
    (5.8,-1) rectangle (7.8,1);
  \node[lbl,align=center,text width=32mm] at (3.4,1.95)
    {every point on each shape has norm exactly $1$};
\end{tikzpicture}
\caption{The set of vectors of norm one, under each of the three norms. The
corners of the $L_1$ diamond sit on the axes, which is the geometric reason
$L_1$ regularisation drives coefficients to exactly zero while $L_2$ merely
shrinks them --- a constrained optimum is far more likely to land on a corner
than on a smooth arc.}
\label{fig:math-unit-balls}
\end{figure}

\begin{engineernote}
Pick the norm from the question, not from habit. $L_2$ is the default because
it is smooth and differentiable everywhere. $L_1$ is the choice when you want
sparsity or robustness to outliers, because it does not square the errors.
$L_\infty$ is the choice when the worst single component is what matters ---
a safety margin, a latency budget, an adversarial perturbation bound.
\end{engineernote}

The \keyterm{unit vector} in the direction of $\vec{x}$ is
$\vec{x}/\lVert\vec{x}\rVert_2$, which has length
one\textsuperscript{\cite[][p.~262]{cherneyLinearAlgebra2013}}. Normalising
throws away magnitude and keeps only direction --- often exactly what you want
when comparing documents of different lengths.

\subsection{The Dot Product and Cosine Similarity}

The \keyterm{dot product} of $\vec{u}$ and $\vec{v}$ multiplies matching
components and sums\textsuperscript{\cite[][p.~89]{cherneyLinearAlgebra2013}}:
\begin{equation}
  \vec{u}\cdot\vec{v} = \sum_{i=1}^{n} u_i v_i
  = \lVert\vec{u}\rVert\,\lVert\vec{v}\rVert\cos\Theta .
  \label{eq:math-dot}
\end{equation}
\where{$u_i,v_i$ are matching components, $n$ the shared dimension, and
$\Theta$ the angle between the two vectors.}
The second equality is the useful one: the dot product secretly measures
\emph{angle}. Two orthogonal vectors have dot product
zero\textsuperscript{\cite[][p.~30]{fischerMaschinellesLernenFuer2024}}.
Rearranging gives \keyterm{cosine similarity}:
\begin{equation}
  \cos\Theta = \frac{\vec{u}\cdot\vec{v}}
                    {\lVert\vec{u}\rVert\,\lVert\vec{v}\rVert} .
  \label{eq:math-cosine}
\end{equation}
\where{$\lVert\vec{u}\rVert,\lVert\vec{v}\rVert$ are the vectors' lengths,
so dividing by them removes magnitude and leaves only direction.}

\begin{workedexample}[Angle between two vectors, by hand]
Let $\vec{u} = (2,1,2)$ and $\vec{v} = (1,2,2)$.

\emph{Step 1 --- dot product.}
$\vec{u}\cdot\vec{v} = 2(1) + 1(2) + 2(2) = 2+2+4 = 8$.

\emph{Step 2 --- lengths.}
$\lVert\vec{u}\rVert = \sqrt{4+1+4} = \sqrt{9} = 3$ and
$\lVert\vec{v}\rVert = \sqrt{1+4+4} = \sqrt{9} = 3$.

\emph{Step 3 --- cosine.}
$\cos\Theta = \dfrac{8}{3\cdot3} = \dfrac{8}{9} = 0.889$.

\emph{Step 4 --- angle.}
$\Theta = \arccos(0.889) = 27.3^\circ$.

\emph{Interpret.} The vectors point in noticeably similar directions. Because
both happen to have length $3$, the normalisation did little work here --- but
if $\vec{v}$ were $(10,20,20)$, ten times longer, the dot product would be $80$
while the cosine would still be $0.889$. That invariance to magnitude is the
entire reason embedding search uses cosine similarity rather than the raw dot
product.
\end{workedexample}

\begin{figure}[htbp]
\centering
\begin{tikzpicture}[scale=1.0,
  v/.style={-{Stealth[length=2.4mm]},very thick},
  lbl/.style={font=\sffamily\scriptsize,color=AlmanachInk!75}]
  \foreach \dx/\ang/\txt/\val in {
      0/12/{nearly aligned}/{$\cos\Theta\approx0.98$},
      4.3/90/{orthogonal}/{$\cos\Theta = 0$},
      8.6/168/{opposed}/{$\cos\Theta\approx-0.98$}}{
    \draw[draw=AlmanachInk!25] (\dx,0) circle (1.15);
    \draw[v,AlmanachTeal] (\dx,0) -- ++(0:1.15);
    \draw[v,AlmanachOchre] (\dx,0) -- ++(\ang:1.15);
    \node[lbl] at (\dx,-1.55) {\txt};
    \node[lbl] at (\dx,-1.95) {\val};}
  \draw[AlmanachRed] (0.55,0) arc (0:12:0.55);
  \node[lbl,color=AlmanachRed] at (0.85,0.18) {$\Theta$};
  \draw[AlmanachRed] (4.85,0) arc (0:90:0.55);
  \draw[AlmanachRed] (9.15,0) arc (0:168:0.55);
\end{tikzpicture}
\caption{Cosine similarity across its whole range. Because the measure depends
only on angle, both vectors can be rescaled arbitrarily without changing the
answer --- which is what makes it the right similarity for embeddings, where
vector magnitude often encodes frequency rather than meaning.}
\label{fig:math-cosine-similarity}
\end{figure}

\begin{cautionbox}[High dimensions do not behave like the page]
Your geometric intuition is trained in two and three dimensions and it will
mislead you in $576$. In high-dimensional space, randomly chosen vectors are
almost always close to orthogonal, and the distances between all pairs of
points become nearly equal --- so ``nearest neighbour'' becomes progressively
less meaningful as dimension grows. This is the \keyterm{concentration of
distances}, one face of the curse of dimensionality. The practical consequence:
a similarity score of $0.3$ may be a very strong match in a $1000$-dimensional
embedding space even though it sounds weak. Always calibrate similarity
thresholds against the actual distribution of scores in your data rather than
against intuition.
\end{cautionbox}

\begin{checkpoint}
Two document embeddings have cosine similarity $0.95$ but Euclidean distance
$40$. Explain how both can be true at once, and state which measure you would
report to a user searching for ``documents about the same topic''.
\end{checkpoint}

% ===========================================================================
\section{Matrices: Data as Transformation}
% ===========================================================================

A matrix $\mathbf{A}\in\mathbb{R}^{m\times n}$ has $m$ rows and $n$ columns.
The product $\mathbf{C} = \mathbf{A}\mathbf{B}$ is defined when
$\mathbf{A}\in\mathbb{R}^{m\times k}$ and $\mathbf{B}\in\mathbb{R}^{k\times n}$,
with $C_{ij} = \sum_{l=1}^{k} A_{il}B_{lj}$\textsuperscript{\cite{cherneyLinearAlgebra2013},
\cite[][ch.~2]{strangIntroductionLinearAlgebra2016}}.

\begin{plainlanguage}
Do not think of a matrix as a table of numbers. Think of it as a \emph{machine}
with a fixed number of input slots and output slots. Feed it a vector of the
right length and it returns a different vector. A neural network layer is
exactly this machine plus one bend. Matrix multiplication is what you get when
you bolt two machines together and ask what the combined machine does.
\end{plainlanguage}

\begin{figure}[htbp]
\centering
\begin{tikzpicture}[
  lbl/.style={font=\sffamily\scriptsize,color=AlmanachInk!75},
  dim/.style={font=\ttfamily\scriptsize,color=AlmanachTeal!80!black},
  bx/.style={draw=AlmanachTeal,fill=AlmanachMist,minimum height=11mm,
    font=\sffamily\small}]
  \node[bx,minimum width=17mm] (a) at (0,0) {$\mathbf{A}$};
  \node[bx,minimum width=13mm] (b) at (2.6,0) {$\mathbf{B}$};
  \node[font=\large] at (1.35,0) {$\cdot$};
  \node[font=\large] at (4.0,0) {$=$};
  \node[bx,minimum width=13mm,fill=AlmanachOchre!18,draw=AlmanachOchre]
    (c) at (5.6,0) {$\mathbf{C}$};
  \node[dim,above=1.5mm of a] {(m, \textbf{k})};
  \node[dim,above=1.5mm of b] {(\textbf{k}, n)};
  \node[dim,above=1.5mm of c] {(m, n)};
  \draw[-{Stealth[length=1.8mm]},AlmanachRed,thick]
    ($(a.south)+(6mm,-1mm)$) .. controls +(0,-7mm) and +(0,-7mm) ..
    ($(b.south)+(-4mm,-1mm)$)
    node[midway,below=1mm,lbl,color=AlmanachRed] {must match, and vanishes};
  \draw[-{Stealth[length=1.8mm]},AlmanachTeal,thick]
    ($(a.north west)+(1mm,3mm)$) .. controls +(0,9mm) and +(0,9mm) ..
    ($(c.north west)+(1mm,3mm)$);
  \draw[-{Stealth[length=1.8mm]},AlmanachTeal,thick]
    ($(b.north east)+(-1mm,3mm)$) .. controls +(0,6mm) and +(0,6mm) ..
    ($(c.north east)+(-1mm,3mm)$);
\end{tikzpicture}
\caption{The shape rule for matrix multiplication. The inner dimensions must
agree and then disappear; the outer dimensions survive into the result. Reading
a chain of products right to left with this rule is enough to catch most
dimension bugs before running anything.}
\label{fig:math-matmul-shape}
\end{figure}

\begin{workedexample}[Matrix multiplication, entry by entry]
Let
\[
  \mathbf{A} = \begin{pNiceMatrix} 1 & 2 & 3\\ 4 & 5 & 6 \end{pNiceMatrix}
  \ (2\times3),
  \qquad
  \mathbf{B} = \begin{pNiceMatrix} 7 & 8\\ 9 & 10\\ 11 & 12 \end{pNiceMatrix}
  \ (3\times2).
\]
Inner dimensions are both $3$, so the product exists and is $2\times2$. Each
entry is one row of $\mathbf{A}$ dotted with one column of $\mathbf{B}$:
\[
\begin{aligned}
  C_{11} &= 1(7) + 2(9) + 3(11) = 7+18+33 = 58, &
  C_{12} &= 1(8) + 2(10) + 3(12) = 8+20+36 = 64,\\
  C_{21} &= 4(7) + 5(9) + 6(11) = 28+45+66 = 139, &
  C_{22} &= 4(8) + 5(10) + 6(12) = 32+50+72 = 154.
\end{aligned}
\]
\[
  \mathbf{C} = \begin{pNiceMatrix} 58 & 64\\ 139 & 154 \end{pNiceMatrix}.
\]
Note that $\mathbf{B}\mathbf{A}$ is also defined here --- but it is
$3\times3$, not $2\times2$. Matrix multiplication is \emph{not} commutative,
and in general $\mathbf{A}\mathbf{B}$ and $\mathbf{B}\mathbf{A}$ do not even
have the same shape.
\end{workedexample}

\subsection{Transpose, Trace, Determinant, Inverse}

Four properties do most of the work\textsuperscript{\cite{cherneyLinearAlgebra2013}}:
\begin{description}
  \item[Transpose $\mathbf{A}^{\!\top}$] Rows become columns. Note the
    order reversal: $(\mathbf{A}\mathbf{B})^{\!\top} =
    \mathbf{B}^{\!\top}\mathbf{A}^{\!\top}$.
  \item[Trace $\operatorname{tr}(\mathbf{A})$] The sum of the diagonal, which
    also equals the sum of the eigenvalues.
  \item[Determinant $\det(\mathbf{A})$] The factor by which the transformation
    scales volume. It is zero exactly when the matrix is singular.
  \item[Inverse $\mathbf{A}^{-1}$] The transformation that undoes
    $\mathbf{A}$. It exists exactly when $\det(\mathbf{A})\neq 0$.
\end{description}
For the $2\times2$ case,
\begin{equation}
  \mathbf{A} = \begin{pNiceMatrix} a & b\\ c & d \end{pNiceMatrix}
  \ \Longrightarrow\
  \det(\mathbf{A}) = ad-bc,
  \qquad
  \mathbf{A}^{-1} = \frac{1}{ad-bc}
  \begin{pNiceMatrix} d & -b\\ -c & a \end{pNiceMatrix}.
  \label{eq:math-2x2-inverse}
\end{equation}
\where{$a,b,c,d$ are the four entries of $\mathbf{A}$, read left to right
and top to bottom.}

\begin{figure}[htbp]
\centering
\begin{tikzpicture}[scale=0.92,
  lbl/.style={font=\sffamily\scriptsize,color=AlmanachInk!75},
  v/.style={-{Stealth[length=2mm]},thick}]
  \begin{scope}
    \draw[draw=AlmanachInk!30,step=0.5,very thin] (-0.2,-0.2) grid (2.2,2.2);
    \fill[AlmanachTeal!22,draw=AlmanachTeal,thick]
      (0,0) -- (1,0) -- (1,1) -- (0,1) -- cycle;
    \draw[v,AlmanachTeal] (0,0) -- (1,0);
    \draw[v,AlmanachOchre] (0,0) -- (0,1);
    \node[lbl] at (1.0,-0.65) {unit square, area $1$};
  \end{scope}
  \draw[-{Stealth[length=2.6mm]},very thick,AlmanachInk!60] (2.7,1) -- (4.1,1)
    node[midway,above,lbl] {$\mathbf{A}$};
  \begin{scope}[xshift=48mm]
    \draw[draw=AlmanachInk!30,step=0.5,very thin] (-0.2,-0.2) grid (3.2,2.7);
    \fill[AlmanachTeal!22,draw=AlmanachTeal,thick]
      (0,0) -- (2,0.5) -- (2.5,2) -- (0.5,1.5) -- cycle;
    \draw[v,AlmanachTeal] (0,0) -- (2,0.5);
    \draw[v,AlmanachOchre] (0,0) -- (0.5,1.5);
    \node[lbl] at (1.5,-0.65) {area $=\lvert\det\mathbf{A}\rvert = 2.75$};
  \end{scope}
\end{tikzpicture}
\caption{What the determinant means. The matrix maps the unit square to a
parallelogram, and the determinant is the factor by which area changed. A
determinant of zero means the square was flattened onto a line --- information
was destroyed, which is exactly why such a matrix cannot be inverted.}
\label{fig:math-determinant-area}
\end{figure}

\begin{workedexample}[Inverting a $2\times2$ matrix and checking the result]
Let $\mathbf{A} = \begin{pNiceMatrix} 4 & 7\\ 2 & 6\end{pNiceMatrix}$.

\emph{Determinant.} $\det(\mathbf{A}) = 4(6) - 7(2) = 24-14 = 10 \neq 0$, so the
inverse exists.

\emph{Inverse.} By \cref{eq:math-2x2-inverse},
\[
  \mathbf{A}^{-1} = \frac{1}{10}
  \begin{pNiceMatrix} 6 & -7\\ -2 & 4\end{pNiceMatrix}
  = \begin{pNiceMatrix} 0.6 & -0.7\\ -0.2 & 0.4\end{pNiceMatrix}.
\]

\emph{Always check.} Multiply back:
\[
  \mathbf{A}\mathbf{A}^{-1} =
  \begin{pNiceMatrix}
    4(0.6)+7(-0.2) & 4(-0.7)+7(0.4)\\
    2(0.6)+6(-0.2) & 2(-0.7)+6(0.4)
  \end{pNiceMatrix}
  = \begin{pNiceMatrix} 2.4-1.4 & -2.8+2.8\\ 1.2-1.2 & -1.4+2.4\end{pNiceMatrix}
  = \begin{pNiceMatrix} 1 & 0\\ 0 & 1\end{pNiceMatrix}. \checkmark
\]
Multiplying back is a two-minute check that catches sign errors, which are by
far the most common mistake in hand-computed inverses.
\end{workedexample}

\subsection{Eigenvalues and Eigenvectors}

A non-zero $\vec{v}$ is an \keyterm{eigenvector} of $\mathbf{A}$ with
\keyterm{eigenvalue} $\lambda$ when $\mathbf{A}\vec{v} =
\lambda\vec{v}$\textsuperscript{\cite{cherneyLinearAlgebra2013}} --- the matrix
scales $\vec{v}$ without rotating it. Eigenvalues solve the characteristic
equation:
\begin{equation}
  \mathbf{A}\vec{v} = \lambda\vec{v}
  \ \Longrightarrow\
  (\mathbf{A}-\lambda\mathbf{I})\vec{v} = \vec{0}
  \ \Longrightarrow\
  \det(\mathbf{A}-\lambda\mathbf{I}) = 0 .
  \label{eq:math-eigen}
\end{equation}
\where{$\lambda$ is a scalar (the eigenvalue), $\vec{v}\neq\vec{0}$ the
corresponding eigenvector, and $\mathbf{I}$ the identity matrix of the same
size as $\mathbf{A}$.}

\begin{plainlanguage}
Most vectors get knocked off course when a matrix acts on them. A few special
directions survive: the matrix stretches or squashes them but leaves them
pointing the same way. Those are the eigenvectors, and the stretch factor is
the eigenvalue. They are the transformation's own natural axes --- the
directions in which it is simplest to describe.
\end{plainlanguage}

\begin{workedexample}[Eigenvalues of a $2\times2$ matrix, with two checks]
Let $\mathbf{A} = \begin{pNiceMatrix} 2 & 1\\ 1 & 2\end{pNiceMatrix}$.

\emph{Characteristic polynomial.}
\[
  \det\!\begin{pNiceMatrix} 2-\lambda & 1\\ 1 & 2-\lambda\end{pNiceMatrix}
  = (2-\lambda)^2 - 1 = \lambda^2 - 4\lambda + 3 = 0 .
\]
Factorising, $(\lambda-1)(\lambda-3) = 0$, so $\lambda_1 = 1$ and
$\lambda_2 = 3$.

\emph{Check 1 --- trace.} $\operatorname{tr}(\mathbf{A}) = 2+2 = 4$ and
$\lambda_1+\lambda_2 = 1+3 = 4$. \checkmark

\emph{Check 2 --- determinant.} $\det(\mathbf{A}) = 4-1 = 3$ and
$\lambda_1\lambda_2 = 1\times3 = 3$. \checkmark

\emph{Eigenvectors.} For $\lambda=3$: $(2-3)v_1 + v_2 = 0 \Rightarrow v_2 = v_1$,
so $\vec{v} = (1,1)$. For $\lambda=1$: $(2-1)v_1 + v_2 = 0 \Rightarrow
v_2 = -v_1$, so $\vec{v} = (1,-1)$.

\emph{Verify directly.}
$\mathbf{A}(1,1) = (2+1,\,1+2) = (3,3) = 3\,(1,1)$. \checkmark

Those two checks --- trace equals the sum, determinant equals the product ---
cost seconds and catch nearly every algebra slip.
\end{workedexample}

\begin{figure}[htbp]
\centering
\begin{tikzpicture}[scale=0.95,
  ax/.style={->,draw=AlmanachInk!45,>={Stealth[length=1.6mm]}},
  lbl/.style={font=\sffamily\scriptsize,color=AlmanachInk!75},
  before/.style={-{Stealth[length=2mm]},thick,draw=AlmanachInk!55},
  after/.style={-{Stealth[length=2.4mm]},very thick,draw=AlmanachOchre}]
  \draw[ax] (-2.4,0) -- (2.6,0); \draw[ax] (0,-2.0) -- (0,2.4);
  \draw[AlmanachTeal!45,dashed,thick] (-2.2,-2.2) -- (2.2,2.2);
  \draw[AlmanachTeal!45,dashed,thick] (-1.8,1.8) -- (1.8,-1.8);
  \draw[before] (0,0) -- (1,1);
  \draw[after]  (0,0) -- (2.05,2.05);
  \node[lbl,color=AlmanachOchre,anchor=west] at (1.55,2.25)
    {$\lambda_2=3$: stretched $3\times$};
  \draw[before] (0,0) -- (1,-1);
  \draw[after]  (0,0) -- (1,-1);
  \node[lbl,color=AlmanachOchre,anchor=west] at (1.1,-1.15)
    {$\lambda_1=1$: unchanged};
  \draw[before] (0,0) -- (0,1.4);
  \draw[after]  (0,0) -- (0.7,1.4);
  \node[lbl,anchor=west] at (0.75,1.15) {a non-eigenvector: rotated};
\end{tikzpicture}
\caption{Eigenvectors of $\mathbf{A}=\left(\begin{smallmatrix}2&1\\1&2
\end{smallmatrix}\right)$. Vectors along the two dashed axes keep their
direction and are only scaled; every other vector is rotated as well. Finding
those axes is what principal component analysis does to a covariance matrix.}
\label{fig:math-eigenvectors}
\end{figure}

\begin{engineernote}
This is where principal component analysis comes from. Compute the covariance
matrix of your data, find its eigenvectors, and you have the directions along
which the data actually varies; the eigenvalues tell you how much variance each
direction explains. ``Keep the top $k$ components'' means ``keep the $k$
eigenvectors with the largest eigenvalues''. Nothing more mysterious than
that.
\end{engineernote}

\begin{cautionbox}[Never actually invert a matrix]
The formula $\hat{\mathbf{w}} = (\mathbf{X}^{\!\top}\mathbf{X})^{-1}
\mathbf{X}^{\!\top}\mathbf{y}$ appears in every textbook, and you should almost
never implement it that way. Forming and inverting
$\mathbf{X}^{\!\top}\mathbf{X}$ \emph{squares} the condition number of the
problem, so it destroys roughly twice as many digits of accuracy as necessary.
Solve the linear system directly instead --- via QR or SVD, which is what
\texttt{lstsq} does in every numerical
library\textsuperscript{\cite[][ch.~20]{highamAccuracyStabilityNumerical2002}}.
\Cref{sec:math-numerics} shows what ill-conditioning does with real numbers.
\end{cautionbox}

\begin{checkpoint}
A matrix has eigenvalues $\{5, 2, 0\}$. State its determinant, its trace,
whether it is invertible, and what happened geometrically to the volume of a
unit cube passed through it.
\end{checkpoint}

% ===========================================================================
\section{Calculus: Slopes, and Why Machines Follow Them}
% ===========================================================================

\subsection{The Derivative}

The \keyterm{derivative} of $y = f(x)$ is the limit of the rise over the
run\textsuperscript{\cite[][p.~89]{yangMathematicsCivilEngineers2018},
\cite[][p.~393]{bronsteinTaschenbuchMathematik2001}}. The three notations
$f'(x)$, $dy/dx$ and $df(x)/dx$ are interchangeable:
\begin{equation}
  f'(x) = \frac{dy}{dx} = \lim_{\Delta x\to0}
  \frac{f(x+\Delta x) - f(x)}{\Delta x}.
  \label{eq:math-derivative}
\end{equation}
\where{$\Delta x$ is a small increment in the input, and the quotient is
the average rate of change over that increment; the limit shrinks it to a
point.}

\begin{plainlanguage}
The derivative answers one question: \emph{if I nudge the input a little, how
much does the output move, and in which direction?} That is the only question
machine learning ever asks of calculus. A positive derivative says ``increasing
the input increases the output''. Training a model means asking that question
about millions of parameters at once and then stepping every one of them
slightly downhill.
\end{plainlanguage}

\begin{figure}[htbp]
\centering
\begin{tikzpicture}
  \begin{axis}[
      width=10cm, height=6.4cm,
      axis lines=middle,
      axis line style={thick,-Stealth,draw=AlmanachInk!60},
      xlabel={$x$}, ylabel={$y=f(x)$},
      xmin=0, xmax=3, ymin=-0.5, ymax=5,
      samples=100, domain=0:2.2,
      legend cell align=left,
      legend style={font=\scriptsize,at={(1.02,1.05)},anchor=north east,
        cells={anchor=west},draw=none,fill=none},
      tick label style={font=\scriptsize},
      label style={font=\small\sffamily},
      ytick distance=1, minor y tick num=0,
      grid=both,
      major grid style={draw=AlmanachInk!12},
      minor grid style={draw=AlmanachInk!6},
      clip=false,
      ylabel style={at={(ticklabel cs:1.0)},anchor=south},
      xlabel style={at={(1.0,0.1)},anchor=south}]
    \addplot[very thick,AlmanachTeal,smooth] {x^2};
    \addlegendentry{$y = x^2$}
    \addplot[thick,AlmanachRed,dashed,domain=0.5:2.5] {2*x - 1};
    \addlegendentry{tangent $y = 2x-1$ at $P$}
    \addplot[only marks,mark=*,mark size=2pt,AlmanachInk]
      coordinates {(1,1)};
    \node[above left,font=\footnotesize] at (axis cs:1,1) {$P$};
    \draw[thick,AlmanachInk!55] (axis cs:1,1) -- (axis cs:1,0);
    \draw[thick,AlmanachInk!55] (axis cs:1,1) -- (axis cs:2,1);
    \draw[thick,AlmanachInk!55] (axis cs:2,0) -- (axis cs:2,4);
    \addplot[only marks,mark=*,mark size=1.5pt,AlmanachInk!70]
      coordinates {(2,1) (2,3) (2,4)};
    \draw[thick,AlmanachOchre,decorate,
      decoration={brace,amplitude=4pt,mirror}]
      (axis cs:2.1,1) -- (axis cs:2.1,3);
    \node[right,font=\footnotesize,color=AlmanachOchre]
      at (axis cs:2.15,2) {$dy$};
    \draw[thick,AlmanachRed,decorate,
      decoration={brace,amplitude=4pt,mirror}]
      (axis cs:2.4,1) -- (axis cs:2.4,4);
    \node[right,font=\footnotesize,color=AlmanachRed]
      at (axis cs:2.45,2.5) {$\Delta y$};
    \draw[thick,AlmanachInk!50,decorate,
      decoration={brace,amplitude=4pt,mirror}]
      (axis cs:1,-0.55) -- (axis cs:2,-0.55);
    \node[below,font=\footnotesize,color=AlmanachInk!70]
      at (axis cs:1.5,-0.7) {$\Delta x = dx$};
    \path (1,0) coordinate (A) -- (0.5,0) coordinate (B) --
      (1,1) coordinate (C)
      pic ["$\alpha$",draw=AlmanachInk!70,-latex,angle radius=25pt] {angle};
  \end{axis}
\end{tikzpicture}
\caption{The derivative as the slope of the tangent. The gap between $dy$ (what
the tangent predicts) and $\Delta y$ (what the function actually does) is the
error of the linear approximation. It shrinks quadratically as $\Delta x\to0$,
which is why small steps are safe and large steps are not --- the same fact
that governs the learning rate.}
\label{fig:derivative}
\end{figure}

\subsection{The Chain Rule}

Of all the rules, one matters more than the rest combined. If a function is a
composition, $y = f(g(x))$, then
\begin{equation}
  \frac{dy}{dx} = \frac{dy}{du}\cdot\frac{du}{dx},
  \qquad u = g(x).
  \label{eq:math-chain-rule}
\end{equation}
\where{$u$ is the intermediate value produced by the inner function, so
$dy/du$ is the outer rate of change and $du/dx$ the inner one.}

\begin{engineernote}
\Cref{eq:math-chain-rule} \emph{is} backpropagation. A neural network is a
composition of a hundred functions; training it means applying the chain rule a
hundred times and multiplying the results. Everything in
\cref{ch:deep-learning} about vanishing and exploding gradients is a statement
about what happens to a product of many such factors. If you learn one formula
from this chapter, learn this one.
\end{engineernote}

\begin{table}[htbp]
\centering
\caption{The differentiation rules used in this book. The chain rule is listed
last because everything else is a special case you will apply inside it.}
\label{tab:math-diff-rules}
\begin{tblr}{
  width=\textwidth,
  colspec={Q[l,0.16\textwidth] Q[l,0.30\textwidth] X[l]},
  row{1}={font=\bfseries,bg=AlmanachMist},
  hlines,vlines}
Rule & Statement & Where it appears \\
Power & $\dfrac{d}{dx}x^{n} = nx^{n-1}$ & Squared-error losses \\
Exponential & $\dfrac{d}{dx}e^{x} = e^{x}$ & Softmax, sigmoid, likelihoods \\
Logarithm & $\dfrac{d}{dx}\ln x = \dfrac{1}{x}$ &
  Cross-entropy, log-likelihood \\
Product & $(uv)' = u'v + uv'$ & Gating in LSTM and attention \\
Quotient & $\left(\dfrac{u}{v}\right)' =
  \dfrac{u'v-uv'}{v^{2}}$ & Deriving the softmax gradient \\
Chain & $\dfrac{dy}{dx} = \dfrac{dy}{du}\dfrac{du}{dx}$ &
  Every layer of every network \\
\end{tblr}
\end{table}

\begin{workedexample}[The chain rule, twice]
\emph{(a) A pure calculus case.} Differentiate $f(x) = (3x^2+1)^5$.

Set $u = 3x^2+1$, so $f = u^5$. Then $df/du = 5u^4$ and $du/dx = 6x$, giving
\[
  \frac{df}{dx} = 5u^4\cdot 6x = 30x\,(3x^2+1)^4 .
\]
At $x=1$: $u = 4$, so $df/dx = 30(1)(4^4) = 30\times256 = 7680$.

\emph{Sanity check by finite difference.} $f(1) = 4^5 = 1024$ and
$f(1.001) = (3.006003+1)^5 = 4.006003^5 = 1031.69$, so
\[
  \frac{1031.69-1024}{0.001} = 7690 \approx 7680 . \checkmark
\]
The small discrepancy is the linear-approximation error of
\cref{fig:derivative}, not an algebra mistake.

\emph{(b) The case you will actually meet.} A single sigmoid unit with loss
$\mathcal{L} = \tfrac12(\sigma(z)-y)^2$ where $z = wx+b$ and
$\sigma(z) = 1/(1+e^{-z})$. Chaining three links,
\[
  \frac{\partial\mathcal{L}}{\partial w}
  = \underbrace{(\sigma(z)-y)}_{\partial\mathcal{L}/\partial\sigma}
  \cdot\underbrace{\sigma(z)\bigl(1-\sigma(z)\bigr)}_{\partial\sigma/\partial z}
  \cdot\underbrace{x}_{\partial z/\partial w}.
\]
With $w=0.5$, $x=2$, $b=0$, $y=1$: $z = 1$, $\sigma(1) = 0.7311$, so
\[
  \frac{\partial\mathcal{L}}{\partial w}
  = (0.7311-1)(0.7311)(0.2689)(2)
  = (-0.2689)(0.1966)(2) = -0.1057 .
\]
The sign is negative, so increasing $w$ \emph{decreases} the loss --- gradient
descent will therefore increase $w$, moving $\sigma(z)$ closer to the target
$y=1$. Exactly what it should do.
\end{workedexample}

\subsection{Partial Derivatives and the Gradient}

For $f(x_1,\dots,x_n)$, the \keyterm{partial derivative} $\partial f/\partial
x_i$ measures the change with respect to $x_i$ while everything else is held
fixed\textsuperscript{\cite{bronsteinTaschenbuchMathematik2001}}. Collecting all
of them gives the \keyterm{gradient}, a vector:
\begin{equation}
  \nabla f =
  \begin{pNiceMatrix}
    \dfrac{\partial f}{\partial x_1}\\[6pt] \vdots\\[2pt]
    \dfrac{\partial f}{\partial x_n}
  \end{pNiceMatrix},
  \qquad
  \text{and gradient descent steps }
  \vec{x} \leftarrow \vec{x} - \eta\,\nabla f .
  \label{eq:math-gradient}
\end{equation}
\where{each entry is the partial derivative with respect to one input, so
$\nabla f$ has the same shape as $\vec{x}$; $\eta>0$ is the step size, called
the learning rate.}
The gradient points in the direction of \emph{steepest ascent}; descent
therefore moves against it.

\begin{workedexample}[A gradient computed by hand]
Let $f(x,y) = x^2y + 3y^2$.
\[
  \frac{\partial f}{\partial x} = 2xy
  \quad(\text{treat }y\text{ as a constant}),
  \qquad
  \frac{\partial f}{\partial y} = x^2 + 6y
  \quad(\text{treat }x\text{ as a constant}).
\]
At the point $(2,1)$: $\nabla f = (2\cdot2\cdot1,\ 4 + 6) = (4,\,10)$.

\emph{Interpret.} The steepest uphill direction from $(2,1)$ is
$(4,10)$ --- mostly in $y$. Moving one small unit in $y$ raises $f$ about two
and a half times as much as moving one unit in $x$. To \emph{minimise} $f$ you
step along $(-4,-10)$.
\end{workedexample}

\begin{figure}[htbp]
\centering
\begin{tikzpicture}[scale=1.0,
  lbl/.style={font=\sffamily\scriptsize,color=AlmanachInk!75},
  ax/.style={->,draw=AlmanachInk!50,>={Stealth[length=1.6mm]}}]
  \draw[ax] (-2.6,0) -- (2.8,0) node[right,lbl] {$x$};
  \draw[ax] (0,-2.2) -- (0,2.5) node[above,lbl] {$y$};
  \foreach \r/\op in {0.6/18, 1.2/32, 1.8/46, 2.4/60}{
    \draw[draw=AlmanachTeal!\op,thick] (0,0) ellipse ({\r} and {\r/1.4142});}
  \node[lbl,anchor=west] at (2.5,1.75) {contours of $f=x^2+2y^2$};
  \fill[AlmanachInk] (1.2,0.85) circle (1.2pt);
  \draw[-{Stealth[length=2.2mm]},very thick,AlmanachOchre]
    (1.2,0.85) -- (1.92,1.71);
  \draw[-{Stealth[length=2.2mm]},very thick,AlmanachRed]
    (1.2,0.85) -- (0.48,-0.01);
  \node[lbl,color=AlmanachOchre,anchor=west] at (1.6,1.95) {$+\nabla f$ (uphill)};
  \node[lbl,color=AlmanachRed,anchor=north west] at (0.15,-0.15)
    {$-\nabla f$ (the descent step)};
  \draw[AlmanachInk!45,dashed] (0.35,1.24) -- (2.05,0.46);
  \node[lbl,anchor=west] at (2.05,0.42) {tangent to the contour};
\end{tikzpicture}
\caption{Contours of $f(x,y)=x^2+2y^2$, where each ring is a set of equal
height. The gradient is always perpendicular to the contour through a point and
points uphill; gradient descent follows the opposite arrow. The rings are
elliptical rather than circular because $y$ is penalised twice as heavily,
which is what makes plain descent zig-zag on badly scaled problems.}
\label{fig:math-gradient-contours}
\end{figure}

\subsection{Gradient Descent, and the Learning Rate That Breaks It}

\begin{workedexample}[Three steps of gradient descent, then divergence]
Minimise $f(x) = x^2 - 4x + 5$. Its derivative is $f'(x) = 2x-4$, so the true
minimum is at $x^\star = 2$ with $f(2) = 1$.

\emph{With $\eta = 0.3$, starting at $x_0 = 0$:}

\smallskip
\begin{tblr}{
  width=0.82\linewidth,
  colspec={Q[c,0.10\linewidth] X[c] X[c] X[c]},
  row{1}={font=\bfseries,bg=AlmanachMist},
  hlines,vlines}
Step & $x$ & $f'(x)=2x-4$ & $f(x)$ \\
0 & $0.000$ & $-4.00$ & $5.000$ \\
1 & $0-0.3(-4.00)=1.200$ & $-1.60$ & $1.640$ \\
2 & $1.2-0.3(-1.60)=1.680$ & $-0.64$ & $1.102$ \\
3 & $1.68-0.3(-0.64)=1.872$ & $-0.26$ & $1.016$ \\
\end{tblr}
\smallskip

Each step closes $40\,\%$ of the remaining distance to $x^\star=2$. Smooth
convergence.

\emph{Now with $\eta = 1.1$, same start:}
\[
  x_1 = 0 - 1.1(-4) = 4.4,\quad
  x_2 = 4.4 - 1.1(4.8) = -0.88,\quad
  x_3 = -0.88 - 1.1(-5.76) = 5.456,\ \dots
\]
The iterate oscillates with growing amplitude and never returns.

\emph{Why, exactly.} For a quadratic with $f''=2$, the update multiplies the
distance from the optimum by $\lvert 1-\eta f''\rvert = \lvert1-2\eta\rvert$.
That factor is below one only when
\[
  0 < \eta < \frac{2}{f''} = \frac{2}{2} = 1 .
\]
So $\eta = 0.3$ converges and $\eta = 1.1$ diverges, and the boundary is
exactly $\eta = 1$. This is the cleanest available statement of why the
learning rate is the hyperparameter that matters most: there is a hard
threshold, set by the curvature, above which training cannot work at all.
\end{workedexample}

\begin{figure}[htbp]
\centering
\begin{tikzpicture}
\begin{axis}[almanach axis,
  height=46mm, domain=-1.4:5.6, samples=120,
  xlabel={$x$}, ylabel={$f(x)$},
  ymin=0, ymax=13,
  legend pos=north west]
  \addplot[AlmanachInk!45,thick] {x^2-4*x+5}; \addlegendentry{$f(x)=x^2-4x+5$}
  \addplot[AlmanachTeal,thick,mark=*,mark size=1.4pt,only marks]
    coordinates {(0,5)(1.2,1.64)(1.68,1.102)(1.872,1.016)};
    \addlegendentry{$\eta=0.3$ converges}
  \addplot[AlmanachRed,thick,mark=square*,mark size=1.4pt,only marks]
    coordinates {(0,5)(4.4,6.76)(-0.88,9.294)(5.456,12.945)};
    \addlegendentry{$\eta=1.1$ diverges}
  \addplot[AlmanachTeal,dashed,thick,mark=none]
    coordinates {(0,5)(1.2,1.64)(1.68,1.102)(1.872,1.016)};
  \addplot[AlmanachRed,dashed,thick,mark=none]
    coordinates {(0,5)(4.4,6.76)(-0.88,9.294)(5.456,12.945)};
\end{axis}
\end{tikzpicture}
\caption{The same starting point under two learning rates. Below the stability
threshold $\eta<2/f''$ the iterates walk into the basin; above it they
alternate sides while climbing. Nothing about the model, the data, or the loss
changed --- only the step size.}
\label{fig:math-gradient-descent}
\end{figure}

\subsection{Second Derivatives, Curvature, and Convexity}

The \keyterm{Hessian} collects all second partial derivatives,
$H_{ij} = \partial^2 f/\partial x_i\partial x_j$. It describes curvature. A
function is \keyterm{convex} when its Hessian is positive semi-definite
everywhere, which means every local minimum is also
global\textsuperscript{\cite[][ch.~3]{boydConvexOptimization2004}}.

\begin{figure}[htbp]
\centering
\begin{tikzpicture}
\begin{axis}[almanach axis,height=40mm,width=0.44\linewidth,
  domain=-2.4:2.4,samples=120,xlabel={$x$},ylabel={$f(x)$},
  ymin=-0.4,ymax=6,title={\small\sffamily convex},
  title style={color=AlmanachTeal}]
  \addplot[AlmanachTeal,very thick] {x^2};
  \addplot[only marks,mark=*,mark size=1.6pt,AlmanachOchre]
    coordinates {(0,0)};
  \node[font=\scriptsize\sffamily,anchor=south] at (axis cs:0,0.35)
    {one minimum};
\end{axis}
\end{tikzpicture}
\hfill
\begin{tikzpicture}
\begin{axis}[almanach axis,height=40mm,width=0.44\linewidth,
  domain=-2.4:2.4,samples=160,xlabel={$x$},ylabel={$f(x)$},
  ymin=-2.6,ymax=4.2,title={\small\sffamily non-convex},
  title style={color=AlmanachOchre}]
  \addplot[AlmanachOchre,very thick] {x^4-4*x^2+0.6*x+1};
  \addplot[only marks,mark=*,mark size=1.6pt,AlmanachRed]
    coordinates {(-1.45,-3.2) (1.38,-3.55)};
  \node[font=\scriptsize\sffamily,anchor=south] at (axis cs:-1.45,-2.1)
    {local};
  \node[font=\scriptsize\sffamily,anchor=south] at (axis cs:1.38,-2.4)
    {global};
\end{axis}
\end{tikzpicture}
\caption{Why convexity is worth so much. On the left, any downhill path reaches
the answer and ``I found a minimum'' means ``I found the minimum''. On the
right --- the situation for every neural network --- where you start decides
where you stop, and no algorithm can certify that a better minimum does not
exist elsewhere.}
\label{fig:math-convexity}
\end{figure}

\begin{cautionbox}[Deep learning is not convex, and that is fine]
Linear and logistic regression are convex, so their optima are unique and
reproducible. Neural networks are emphatically not, which sounds alarming and
mostly is not: in high dimensions, bad local minima turn out to be rare, and
most stationary points that trap optimisers are saddle points which noise
escapes. What you \emph{do} lose is guarantees. Two training runs of the same
architecture on the same data with different seeds will land on different
parameters --- so always report results across several seeds, never one.
\end{cautionbox}

\subsection{Integration}

The \keyterm{indefinite integral} is the anti-derivative: if $F'(x) = f(x)$
then $\int f(x)\,dx = F(x)+C$. The \keyterm{definite integral} gives the signed
area under the curve\textsuperscript{\cite{yangMathematicsCivilEngineers2018}}:
\begin{equation}
  \int_a^b f(x)\,dx = F(b)-F(a).
  \label{eq:math-integral}
\end{equation}
\where{$F$ is any antiderivative of $f$, and $[a,b]$ the interval; the result is
the signed area, so regions below the axis subtract.}
The \keyterm{fundamental theorem of calculus} ties the two operations together,
$\frac{d}{dx}\int_a^x f(t)\,dt = f(x)$.

\begin{engineernote}
Integration matters here for exactly one reason: probability of a continuous
variable is \emph{area under a density}, never the height of the curve. When
\cref{fig:normal-histogram-density} shows a bell curve reaching $0.399$ at its
peak, that is a density value, not a probability --- densities may exceed $1$.
Every expectation, every $p$-value, and every confidence interval in this book
is an integral of a density.
\end{engineernote}

\begin{practicallab}[Check a derivative the way a numerical library does]
\textbf{Goal.} Convince yourself that automatic differentiation computes the
same thing your algebra does --- and learn the check that catches a wrong
gradient.

\textbf{Protocol.}
\begin{enumerate}
  \item Pick $f(x) = (3x^2+1)^5$ and differentiate it by hand at $x=1$. You
        should get $7680$.
  \item Compute the central finite difference
        $\bigl(f(x+h)-f(x-h)\bigr)/2h$ for $h = 10^{-2}, 10^{-4}, 10^{-6},
        10^{-8}$.
  \item Plot the absolute error against $h$ on log axes. It falls, reaches a
        minimum, and then \emph{rises} again.
  \item Explain the rise. (It is not a bug in your code --- see
        \cref{sec:math-numerics}.)
  \item Repeat with any autodiff library and confirm it matches your analytic
        answer to machine precision, with no $h$ to choose.
\end{enumerate}

\textbf{Why this matters.} Gradient checking by finite difference is the
standard way to verify a hand-written backward pass. Knowing that the error
curve has a floor tells you which disagreements are real bugs and which are
just arithmetic.
\end{practicallab}

\begin{checkpoint}
You are minimising a loss whose second derivative at the current point is
$f''=50$. Give the largest learning rate that can converge. Then explain why a
network with badly scaled inputs effectively has a very large $f''$ in some
directions, and what that forces you to do to $\eta$.
\end{checkpoint}

% ===========================================================================
\section{Probability: Reasoning Under Uncertainty}
% ===========================================================================

\subsection{Random Variables and Distributions}

A \keyterm{random variable} maps outcomes of a random process to numbers. It is
\keyterm{discrete} when it takes countably many values (a die, a class label)
and \keyterm{continuous} when it takes values on an interval (a temperature, a
latency). Discrete variables have a \keyterm{probability mass function} where
$P(X=x)$ is a genuine probability; continuous variables have a
\keyterm{probability density function} where only areas are
probabilities\textsuperscript{\cite[][ch.~1]{wassermanAllStatistics2004}}.

\begin{cautionbox}[Density is not probability]
For a continuous variable, $P(X = 3.7) = 0$ exactly --- a single point has no
width and therefore no area. Only intervals have probability. A density value
may be greater than $1$ without anything being wrong: a uniform distribution on
$[0,\,0.5]$ has density $2$ everywhere on that interval. If you ever find
yourself reading a probability off the $y$-axis of a density plot, stop.
\end{cautionbox}

\begin{figure}[htbp]
\centering
\begin{tikzpicture}
\begin{axis}[almanach axis,height=38mm,width=0.31\linewidth,
  ybar,bar width=5pt,xlabel={$x$},ylabel={$P(X{=}x)$},
  ymin=0,ymax=0.26,xmin=0.3,xmax=6.7,xtick={1,...,6},
  title={\small\sffamily PMF (discrete)},title style={color=AlmanachTeal}]
  \addplot[fill=AlmanachTeal!35,draw=AlmanachTeal]
    coordinates {(1,0.1667)(2,0.1667)(3,0.1667)(4,0.1667)(5,0.1667)(6,0.1667)};
\end{axis}
\end{tikzpicture}
\hfill
\begin{tikzpicture}
\begin{axis}[almanach axis,height=38mm,width=0.31\linewidth,
  domain=-3.2:3.2,samples=120,xlabel={$x$},ylabel={$f(x)$},
  ymin=0,ymax=0.46,
  title={\small\sffamily PDF (continuous)},title style={color=AlmanachTeal}]
  \addplot[AlmanachTeal,thick,fill=AlmanachTeal!18] {0.39894228*exp(-x^2/2)}
    \closedcycle;
  \node[font=\scriptsize\sffamily,anchor=south] at (axis cs:0,0.05)
    {area $=$ probability};
\end{axis}
\end{tikzpicture}
\hfill
\begin{tikzpicture}
\begin{axis}[almanach axis,height=38mm,width=0.31\linewidth,
  domain=-3.2:3.2,samples=120,xlabel={$x$},ylabel={$F(x)$},
  ymin=0,ymax=1.08,
  title={\small\sffamily CDF},title style={color=AlmanachTeal}]
  \addplot[AlmanachOchre,thick]
    {0.5*(1+tanh(0.7978845608*(x+0.044715*x^3)))};
  \draw[AlmanachInk!35,dashed] (axis cs:-3.2,1) -- (axis cs:3.2,1);
\end{axis}
\end{tikzpicture}
\caption{The three ways to describe one distribution. The mass function gives
probabilities directly; the density must be integrated to give them; the
cumulative function $F(x)=P(X\le x)$ climbs monotonically to one and is what
you actually read quantiles and $p$-values from.}
\label{fig:math-pmf-pdf-cdf}
\end{figure}

\subsection{Expectation and Variance}

The \keyterm{expectation} is the probability-weighted average; the
\keyterm{variance} is the expected squared deviation from it:
\begin{equation}
  \mathbb{E}[X] = \sum_x x\,p(x)
  \quad\text{or}\quad \int x f(x)\,dx,
  \qquad
  \operatorname{Var}(X) = \mathbb{E}\bigl[(X-\mathbb{E}[X])^2\bigr]
  = \mathbb{E}[X^2] - \bigl(\mathbb{E}[X]\bigr)^2 .
  \label{eq:math-expectation}
\end{equation}
\where{$X$ is the random variable, $p(x)$ its probability mass function in
the discrete case and $f(x)$ its density in the continuous one.}

\begin{workedexample}[Expectation and variance of a fair die]
$X$ takes values $1$ to $6$ each with probability $1/6$.

\emph{Expectation.}
$\mathbb{E}[X] = \frac{1+2+3+4+5+6}{6} = \frac{21}{6} = 3.5$.

Note that $3.5$ is not a value the die can show. An expectation is a balance
point, not a prediction.

\emph{Second moment.}
$\mathbb{E}[X^2] = \frac{1+4+9+16+25+36}{6} = \frac{91}{6} = 15.167$.

\emph{Variance} via the shortcut in \cref{eq:math-expectation}:
\[
  \operatorname{Var}(X) = 15.167 - (3.5)^2 = 15.167 - 12.25 = 2.917
  = \tfrac{35}{12},
  \qquad \sigma = \sqrt{2.917} = 1.708 .
\]
The shortcut $\mathbb{E}[X^2]-\mathbb{E}[X]^2$ is worth memorising; it turns a
two-pass calculation into a one-pass one, which is how streaming variance
estimators work.
\end{workedexample}

\subsection{Covariance, Correlation, and a Warning}

\keyterm{Covariance} measures whether two variables move together;
\keyterm{correlation} is covariance rescaled to $[-1,1]$:
\begin{equation}
  \operatorname{Cov}(X,Y) = \mathbb{E}\bigl[(X-\mu_X)(Y-\mu_Y)\bigr],
  \qquad
  \rho_{XY} = \frac{\operatorname{Cov}(X,Y)}{\sigma_X\sigma_Y}.
  \label{eq:math-covariance}
\end{equation}
\where{$\mu_X,\mu_Y$ are the two means and $\sigma_X,\sigma_Y$ the two
standard deviations; dividing by them is what confines $\rho_{XY}$ to
$[-1,1]$.}

\begin{workedexample}[Covariance and correlation from four points]
Let $x = (1,2,3,4)$ and $y = (2,4,5,9)$, so $\bar x = 2.5$ and $\bar y = 5$.

\smallskip
\begin{tblr}{
  width=0.8\linewidth,
  colspec={Q[c]Q[c]Q[c]Q[c]Q[c]},
  row{1}={font=\bfseries,bg=AlmanachMist},
  hlines,vlines}
$x_i$ & $y_i$ & $x_i-\bar x$ & $y_i-\bar y$ & product \\
1 & 2 & $-1.5$ & $-3$ & $4.5$ \\
2 & 4 & $-0.5$ & $-1$ & $0.5$ \\
3 & 5 & $+0.5$ & $\ \ 0$ & $0.0$ \\
4 & 9 & $+1.5$ & $+4$ & $6.0$ \\
\end{tblr}
\smallskip

$\operatorname{Cov} = \frac{4.5+0.5+0+6.0}{4} = \frac{11}{4} = 2.75$.

$\operatorname{Var}(x) = \frac{2.25+0.25+0.25+2.25}{4} = 1.25$ and
$\operatorname{Var}(y) = \frac{9+1+0+16}{4} = 6.5$, so
\[
  \rho = \frac{2.75}{\sqrt{1.25\times6.5}} = \frac{2.75}{2.850} = 0.965 .
\]
Strongly positive. Notice that covariance alone ($2.75$) is uninterpretable
without knowing the scales --- multiply $y$ by $1000$ and the covariance
becomes $2750$ while the correlation stays at $0.965$.
\end{workedexample}

\begin{cautionbox}[Zero correlation does not mean independence]
Correlation detects \emph{linear} association only. Take $X$ uniform on
$[-1,1]$ and $Y = X^2$. Then $Y$ is perfectly determined by $X$ --- as
dependent as two variables can be --- yet $\operatorname{Cov}(X,Y) =
\mathbb{E}[X^3] - \mathbb{E}[X]\mathbb{E}[X^2] = 0-0 = 0$, so $\rho = 0$.
Independence implies zero correlation; the reverse is false. And neither
implies causation: for that you need a causal model and, usually, an
intervention\textsuperscript{\cite{pearlCausalInferenceStatistics2016}}.
\end{cautionbox}

\subsection{Important Distributions}

The \keyterm{normal} (Gaussian) distribution $\mathcal{N}(\mu,\sigma^2)$ is
symmetric about $\mu$, with $\sigma$ setting the
width\textsuperscript{\cite{brucePracticalStatisticsData2017}}. The
\keyterm{68--95--99.7 rule} states that roughly $68\,\%$, $95\,\%$ and
$99.7\,\%$ of the mass lies within one, two and three standard deviations of
the mean.

\factvisual{
  \begin{equation*}
    f(x) = \frac{1}{\sigma\sqrt{2\pi}}
    \exp\!\left(-\frac{(x-\mu)^2}{2\sigma^2}\right)
  \end{equation*}
  \where{$\mu$ is the mean, $\sigma$ the standard deviation, and $x$ the value
  whose density is being evaluated.}
  \begin{itemize}
    \item \textbf{Centre:} mean, median and mode all coincide at $\mu$.
    \item \textbf{Width:} larger $\sigma$ gives a lower, wider bell.
    \item \textbf{Area:} probability is area under the density, never bar
      height.
    \item \textbf{Standardise:} $z=(x-\mu)/\sigma$ converts any value into
      standard-deviation units.
  \end{itemize}
}{
  \begin{tikzpicture}
    \begin{axis}[
      almanach axis,
      height=52mm,
      xmin=-3.4,xmax=3.4,
      ymin=0,ymax=0.46,
      xlabel={$z$-score},
      ylabel={density},
      ytick=\empty,
      xtick={-3,-2,-1,0,1,2,3},
      legend style={at={(0.5,1.02)},anchor=south},
      legend columns=2]
      \addplot+[ybar interval,
        fill=AlmanachTeal!24,
        draw=AlmanachTeal!45,
        mark=none] coordinates {
        (-3.2,0.008)(-2.8,0.015)(-2.4,0.030)
        (-2.0,0.060)(-1.6,0.120)(-1.2,0.205)
        (-0.8,0.300)(-0.4,0.375)(0,0.405)
        (0.4,0.370)(0.8,0.300)(1.2,0.205)
        (1.6,0.120)(2.0,0.060)(2.4,0.030)
        (2.8,0.015)(3.2,0.008)};
      \addlegendentry{sample histogram}
      \addplot[AlmanachOchre,very thick,smooth,
        domain=-3.3:3.3,samples=120]
        {0.39894228*exp(-x^2/2)};
      \addlegendentry{$\mathcal{N}(0,1)$ density}
      \draw[AlmanachInk!45,dashed]
        (axis cs:-1,0) -- (axis cs:-1,0.242);
      \draw[AlmanachInk!45,dashed]
        (axis cs:1,0) -- (axis cs:1,0.242);
      \draw[decorate,decoration={brace,mirror,amplitude=3pt},
        AlmanachTeal]
        (axis cs:-1,-0.012) -- (axis cs:1,-0.012)
        node[midway,below=3pt,font=\scriptsize] {$\approx68\%$};
    \end{axis}
  \end{tikzpicture}
  \visualcaption{Observations form a histogram; the Normal model is a
    continuous density fitted to them. The central interval
    $[-1\sigma,+1\sigma]$ contains about $68\,\%$ of its area.}
    {fig:normal-histogram-density}
}

The \keyterm{binomial} distribution $B(n,p)$ counts successes in $n$
independent trials each with probability
$p$\textsuperscript{\cite{downeyThinkBayes2016}}:
\begin{equation}
  P(X=k) = \binom{n}{k}p^k(1-p)^{n-k},
  \qquad \mu = np,\qquad \sigma^2 = np(1-p).
  \label{eq:math-binomial}
\end{equation}
\where{$n$ is the number of trials, $k$ the number of successes, $p$ the
per-trial success probability, and $\binom{n}{k}$ the number of distinct ways
$k$ successes can occur among $n$ trials.}

\subsection{The Central Limit Theorem}

Sums and averages of many independent contributions tend toward a normal
distribution \emph{whatever the shape of the original distribution}, provided
the variance is finite. This is why the normal distribution is everywhere: it
is not that nature prefers bells, but that averaging manufactures them.

\begin{figure}[htbp]
\centering
\begin{tikzpicture}
\begin{axis}[almanach axis,height=36mm,width=0.31\linewidth,
  ybar,bar width=5.5pt,ymin=0,ymax=0.24,xlabel={one die},
  xmin=0.3,xmax=6.7,xtick={2,4,6},ytick=\empty,
  title={\small\sffamily $n=1$: flat},title style={color=AlmanachTeal}]
  \addplot[fill=AlmanachTeal!35,draw=AlmanachTeal] coordinates
    {(1,.1667)(2,.1667)(3,.1667)(4,.1667)(5,.1667)(6,.1667)};
\end{axis}
\end{tikzpicture}
\hfill
\begin{tikzpicture}
\begin{axis}[almanach axis,height=36mm,width=0.31\linewidth,
  ybar,bar width=3.4pt,ymin=0,ymax=0.20,xlabel={sum of two},
  xmin=1.3,xmax=12.7,xtick={2,7,12},ytick=\empty,
  title={\small\sffamily $n=2$: triangular},title style={color=AlmanachTeal}]
  \addplot[fill=AlmanachTeal!35,draw=AlmanachTeal] coordinates
    {(2,.0278)(3,.0556)(4,.0833)(5,.1111)(6,.1389)(7,.1667)
     (8,.1389)(9,.1111)(10,.0833)(11,.0556)(12,.0278)};
\end{axis}
\end{tikzpicture}
\hfill
\begin{tikzpicture}
\begin{axis}[almanach axis,height=36mm,width=0.31\linewidth,
  ybar,bar width=1.5pt,ymin=0,ymax=0.10,xlabel={sum of five},
  xmin=4,xmax=31,xtick={5,17.5,30},ytick=\empty,
  title={\small\sffamily $n=5$: bell},title style={color=AlmanachTeal}]
  \addplot[fill=AlmanachTeal!35,draw=AlmanachTeal] coordinates
    {(5,.0001)(6,.0006)(7,.0019)(8,.0045)(9,.0090)(10,.0154)(11,.0231)
     (12,.0316)(13,.0401)(14,.0473)(15,.0525)(16,.0548)(17,.0540)
     (18,.0505)(19,.0446)(20,.0373)(21,.0294)(22,.0218)(23,.0150)
     (24,.0096)(25,.0056)(26,.0029)(27,.0013)(28,.0005)(29,.0001)};
  \addplot[AlmanachOchre,thick,smooth,domain=5:30,samples=80]
    {0.0548*exp(-(x-17.5)^2/(2*3.82^2))};
\end{axis}
\end{tikzpicture}
\caption{The central limit theorem visible in three steps. A single die is
uniform; the sum of two is already triangular; the sum of five is
indistinguishable from the fitted normal curve. Nothing about the die changed
--- only how many were added.}
\label{fig:math-clt}
\end{figure}

\subsection{Bayes' Theorem}

\keyterm{Bayes' theorem} is the rule for updating a belief when evidence
arrives:
\begin{equation}
  \underbrace{P(H\mid E)}_{\text{posterior}}
  = \frac{\overbrace{P(E\mid H)}^{\text{likelihood}}\;
          \overbrace{P(H)}^{\text{prior}}}
         {\underbrace{P(E)}_{\text{evidence}}},
  \qquad
  P(E) = P(E\mid H)P(H) + P(E\mid\neg H)P(\neg H).
  \label{eq:math-bayes}
\end{equation}
\where{$H$ is the hypothesis under test and $E$ the observed evidence;
$\neg H$ is the hypothesis being false. The denominator sums the two ways the
evidence could have arisen, which is why it normalises the result.}

\begin{plainlanguage}
Bayes' theorem says: how strongly you should believe something after seeing
evidence depends not only on how convincing the evidence is, but on how likely
the thing was to begin with. Extraordinary claims need extraordinary evidence,
and this is the equation that says exactly how much.
\end{plainlanguage}

\begin{workedexample}[The $99\,\%$ accurate test that is right half the time]
A disease affects $1\,\%$ of the population. A test detects it in $99\,\%$ of
sick people (sensitivity) and correctly clears $99\,\%$ of healthy people
(specificity). You test positive. What is the probability you are ill?

Almost everyone answers $99\,\%$. Work it out.

\emph{Set up.} $P(D) = 0.01$, $P(+\mid D) = 0.99$, $P(+\mid\neg D) = 0.01$.

\emph{Evidence.}
\[
  P(+) = \underbrace{0.99\times0.01}_{\text{true positives}}
       + \underbrace{0.01\times0.99}_{\text{false positives}}
       = 0.0099 + 0.0099 = 0.0198 .
\]

\emph{Posterior.}
\[
  P(D\mid +) = \frac{0.0099}{0.0198} = \boxed{0.50}
\]

\emph{Why.} Look at the two terms in $P(+)$: they are \emph{equal}. In a
population of $10{,}000$ there are $100$ sick people, of whom $99$ test
positive --- and $9{,}900$ healthy people, of whom $99$ also test positive.
The test is excellent, but there are so many more healthy people that false
positives exactly match true ones. A coin flip.

This is not a trick question; it is the everyday arithmetic of screening,
fraud detection, and any classifier deployed on a rare class.
\end{workedexample}

\begin{figure}[htbp]
\centering
\begin{tikzpicture}[
  lbl/.style={font=\sffamily\scriptsize,color=AlmanachInk!80},
  tiny/.style={font=\sffamily\tiny,color=AlmanachInk!70}]
  \def\W{88}
  \fill[AlmanachMist,draw=AlmanachInk!30] (0,0) rectangle (\W mm,34mm);
  \fill[AlmanachOchre!30,draw=AlmanachOchre]
    (0,0) rectangle (0.01*\W mm,34mm);
  \draw[AlmanachInk!30] (0.01*\W mm,0) -- (0.01*\W mm,34mm);
  \fill[AlmanachRed!45,draw=AlmanachRed]
    (0,0) rectangle (0.0099*\W mm,33.66mm);
  \fill[AlmanachRed!45,draw=AlmanachRed]
    (0.01*\W mm,0) rectangle (0.0199*\W mm,34mm);
  \node[tiny,anchor=south west] at (0.5mm,35mm) {sick 1\%};
  \node[tiny,anchor=south west] at (0.022*\W mm,35mm) {healthy 99\%};
  \draw[-{Stealth[length=1.6mm]},AlmanachRed]
    (0.055*\W mm,-7mm) -- (0.014*\W mm,-1mm);
  \node[lbl,anchor=west,color=AlmanachRed,align=left,text width=62mm]
    at (0.06*\W mm,-8mm)
    {The two red slivers are the positive tests: 99 true, 99 false.
     Equal areas $\Rightarrow$ $P(\text{sick}\mid+)=1/2$.};
\end{tikzpicture}
\caption{The base-rate effect drawn to scale. The healthy population is so much
wider that its $1\,\%$ false-positive rate produces exactly as many positive
tests as the sick population's $99\,\%$ true-positive rate. Test quality alone
never determines what a positive result means.}
\label{fig:math-base-rate}
\end{figure}

\begin{engineernote}
This is the mathematics behind every complaint that a fraud detector ``cries
wolf''. If fraud is $0.1\,\%$ of transactions, even a $99.9\,\%$-specific model
generates roughly one false alarm per true one. The fix is never ``make the
model more accurate'' in the abstract --- it is to decide which error is more
costly and move the decision threshold accordingly, then report precision and
recall rather than accuracy. \Cref{ch:evaluation-mlops} takes this apart
properly.
\end{engineernote}

\begin{checkpoint}
Repeat the worked example with a disease prevalence of $10\,\%$ instead of
$1\,\%$, keeping both test rates at $99\,\%$. Compute $P(D\mid+)$. Then state
in one sentence what changed and why screening programmes care so much about
who gets tested.
\end{checkpoint}

\subsection{Bayesian and Frequentist Inference}

Bayesian inference applies \cref{eq:math-bayes} not to a single hypothesis but
to a whole distribution over model parameters
$\theta$\textsuperscript{\cite{downeyThinkBayes2016}}:
\begin{equation}
  P(\theta\mid D) = \frac{P(D\mid\theta)\,P(\theta)}{P(D)},
  \qquad
  P(D) = \int P(D\mid\theta)\,P(\theta)\,d\theta .
  \label{eq:math-bayesian-posterior}
\end{equation}
\where{$\theta$ is the vector of model parameters and $D$ the observed
data; the integral runs over every parameter value the prior allows, which is
what makes it intractable in all but simple models.}
The result is a full posterior distribution rather than a single number, so the
remaining uncertainty is part of the answer instead of a footnote to it. The
integral in the denominator is usually intractable, which is why Markov chain
Monte Carlo and variational inference exist.

\begin{table}[htbp]
\centering
\caption{Two philosophies of probability. Neither is wrong; they answer
different questions, and the practical difference shows up in what you are
allowed to say about the result.}
\label{tab:math-bayes-vs-frequentist}
\begin{tblr}{
  width=\textwidth,
  colspec={Q[l,0.17\textwidth] X[l] X[l]},
  row{1}={font=\bfseries,bg=AlmanachMist},
  hlines,vlines}
Aspect & Frequentist & Bayesian \\
Probability means & Long-run frequency over repeated experiments &
  Degree of belief given what is known \\
Parameters are & Fixed but unknown constants &
  Random variables with a prior distribution \\
Inference yields & Confidence intervals, $p$-values &
  A posterior distribution \\
Data are treated as & One draw from a repeatable process &
  Fixed evidence to condition on \\
Prior knowledge & Not represented & Stated explicitly, and open to challenge \\
You may say & ``$95\,\%$ of such intervals contain the true value'' &
  ``There is a $95\,\%$ probability the value is in this interval'' \\
\end{tblr}
\end{table}

\begin{engineernote}
The last row is the one that trips people up. A frequentist $95\,\%$ confidence
interval does \emph{not} give a $95\,\%$ probability that the parameter is
inside it --- the parameter is a fixed constant, so it either is or is not. The
$95\,\%$ describes the procedure across hypothetical repetitions. Only a
Bayesian credible interval licenses the sentence everyone wants to
say\textsuperscript{\cite[][p.~9]{barberBayesianReasoningMachine2012}}. If you
report an interval, know which kind you computed.
\end{engineernote}

% ===========================================================================
\section{Information Theory: Measuring Surprise}
% ===========================================================================

\keyterm{Entropy} quantifies average uncertainty in a random
variable\textsuperscript{\cite{shannonMathematicalTheoryCommunication},
\cite{mackayInformationTheoryInference2005}}:
\begin{equation}
  H(X) = -\sum_x p(x)\log_2 p(x) \quad[\text{bits}].
  \label{eq:math-entropy}
\end{equation}
\where{$p(x)$ is the probability of outcome $x$ and the sum runs over all
outcomes; using $\log_2$ gives bits, and the natural logarithm gives nats.}

\begin{plainlanguage}
Entropy is the average number of yes/no questions you need to pin down the
answer. A fair coin needs exactly one question, so its entropy is one bit. A
coin that lands heads $99\,\%$ of the time barely needs any question at all,
because you can just guess heads --- so its entropy is near zero. Entropy
measures how surprised you expect to be.
\end{plainlanguage}

\begin{workedexample}[Entropy of fair and biased coins]
\emph{Fair coin}, $p = 0.5$:
\[
  H = -0.5\log_2 0.5 - 0.5\log_2 0.5 = -0.5(-1) - 0.5(-1) = 1.000 \text{ bit.}
\]

\emph{Biased coin}, $p = 0.9$:
\[
  H = -0.9\log_2 0.9 - 0.1\log_2 0.1
    = -0.9(-0.152) - 0.1(-3.322)
    = 0.137 + 0.332 = 0.469 \text{ bits.}
\]

\emph{Certain coin}, $p = 1.0$: $H = 0$ bits --- no uncertainty, no
information in the outcome.

Notice the asymmetry in the biased case: the \emph{rare} outcome contributes
more ($0.332$) than the common one ($0.137$) despite occurring nine times less
often. Rare events carry more information when they happen, which is exactly
why $\log(1/p)$ is the right measure of surprise.
\end{workedexample}

\begin{figure}[htbp]
\centering
\begin{tikzpicture}
\begin{axis}[almanach axis,height=44mm,
  domain=0.001:0.999,samples=200,
  xlabel={$p$ (probability of heads)},
  ylabel={$H(p)$ in bits},
  ymin=0,ymax=1.12,xmin=0,xmax=1]
  \addplot[AlmanachTeal,very thick]
    {-(x*ln(x)+(1-x)*ln(1-x))/ln(2)};
  \addplot[only marks,mark=*,mark size=1.5pt,AlmanachOchre]
    coordinates {(0.5,1) (0.9,0.469)};
  \node[font=\scriptsize\sffamily,anchor=south] at (axis cs:0.5,1.02)
    {fair: 1 bit};
  \node[font=\scriptsize\sffamily,anchor=west] at (axis cs:0.915,0.5)
    {$p{=}0.9$: 0.47};
  \draw[AlmanachInk!35,dashed] (axis cs:0.5,0) -- (axis cs:0.5,1);
\end{axis}
\end{tikzpicture}
\caption{Binary entropy. Uncertainty is maximal at $p=0.5$ and falls to zero at
both extremes, where the outcome is certain. The curve is flat near its peak,
which is why a class balance of $45/55$ is barely easier to predict than
$50/50$, while $1/99$ is a completely different problem.}
\label{fig:math-binary-entropy}
\end{figure}

\subsection{Cross-Entropy and KL Divergence}

\keyterm{Cross-entropy} measures the cost of encoding samples from the true
distribution $p$ using a model $q$; \keyterm{KL divergence} is the excess cost
over the best possible\textsuperscript{\cite{kullbackInformationSufficiency1951},
\cite[][ch.~2]{coverElementsInformationTheory2006}}:
\begin{equation}
  H(p,q) = -\sum_x p(x)\log q(x),
  \qquad
  D_{\mathrm{KL}}(p\,\|\,q) = \sum_x p(x)\log\frac{p(x)}{q(x)}
  = H(p,q) - H(p).
  \label{eq:math-kl}
\end{equation}
\where{$p$ is the true distribution and $q$ the model's; both sums run over
all outcomes. Note the asymmetry: $p$ weights the terms in both expressions,
so swapping $p$ and $q$ changes the answer.}
Since $H(p)$ does not depend on the model, \emph{minimising cross-entropy is
identical to minimising KL divergence} --- which is why cross-entropy is the
default classification loss.

\begin{workedexample}[KL divergence is not symmetric]
Let $p = (0.5,\,0.5)$ be the truth and $q = (0.9,\,0.1)$ the model.
\[
  D_{\mathrm{KL}}(p\,\|\,q)
  = 0.5\log_2\frac{0.5}{0.9} + 0.5\log_2\frac{0.5}{0.1}
  = 0.5(-0.848) + 0.5(2.322)
  = 0.737 \text{ bits.}
\]
Now swap the roles:
\[
  D_{\mathrm{KL}}(q\,\|\,p)
  = 0.9\log_2\frac{0.9}{0.5} + 0.1\log_2\frac{0.1}{0.5}
  = 0.9(0.848) + 0.1(-2.322)
  = 0.531 \text{ bits.}
\]
$0.737 \neq 0.531$. KL is \emph{not} a distance --- it fails symmetry and the
triangle inequality. The direction encodes a real modelling choice:
$D_{\mathrm{KL}}(p\|q)$ punishes the model heavily for assigning near-zero
probability to something that actually happens (mode-covering), while
$D_{\mathrm{KL}}(q\|p)$ punishes it for putting mass where there is none
(mode-seeking). Variational inference and the VAE objective in
\cref{ch:deep-learning} use the second form, which is part of why VAE samples
can look blurry.
\end{workedexample}

\begin{checkpoint}
A classifier outputs $q = (0.7, 0.2, 0.1)$ and the true label is class 1, so
$p = (1,0,0)$. Compute the cross-entropy loss in bits. Then explain why
$D_{\mathrm{KL}}(p\|q)$ equals the cross-entropy exactly in this case.
\end{checkpoint}

% ===========================================================================
\section{Statistics: From Sample to Claim}
% ===========================================================================

\subsection{Describing a Sample}

\keyterm{Central tendency} is summarised by the mean (the balance point), the
median (the middle value) and the mode (the most frequent value).
\keyterm{Spread} is summarised by the variance and standard deviation, or by
the interquartile range.

\begin{workedexample}[Why the median survives an outlier and the mean does not]
Five salaries, in thousands of euro: $40,\ 42,\ 45,\ 47,\ 50$.
\[
  \text{mean} = \frac{224}{5} = 44.8,
  \qquad \text{median} = 45 .
\]
Now the company hires a chief executive on $1000$:
\[
  \text{mean} = \frac{1224}{6} = 204.0,
  \qquad \text{median} = \frac{45+47}{2} = 46 .
\]
The mean rose by a factor of $4.6$; the median moved by $1$. A single
observation moved the mean past every actual salary in the room. When a
distribution is skewed or contaminated, the median is not a rougher summary ---
it is a more honest one.
\end{workedexample}

\subsection{Maximum Likelihood Estimation}

\keyterm{Maximum likelihood} picks the parameter that makes the observed data
most probable\textsuperscript{\cite[][ch.~9]{wassermanAllStatistics2004}}.

\begin{workedexample}[MLE for a coin, derived]
You flip a coin $n=10$ times and see $k=7$ heads. Which $p$ makes that most
likely?

\emph{Likelihood.} $L(p) = \binom{10}{7}p^7(1-p)^3$.

\emph{Take logs} (the maximum is in the same place, and sums are easier than
products):
\[
  \ell(p) = \text{const} + 7\ln p + 3\ln(1-p).
\]

\emph{Differentiate and set to zero:}
\[
  \frac{d\ell}{dp} = \frac{7}{p} - \frac{3}{1-p} = 0
  \ \Longrightarrow\ 7(1-p) = 3p
  \ \Longrightarrow\ 7 = 10p
  \ \Longrightarrow\ \hat p = 0.7 .
\]
The estimate is just the observed frequency --- which is reassuring, and also a
warning: with $n=10$ this estimate is extremely uncertain. A $95\,\%$ interval
runs roughly from $0.35$ to $0.93$. The MLE gives you a number; it does not
give you the right to trust it.
\end{workedexample}

For linear regression, minimising squared error is maximum likelihood under
Gaussian noise, and gives the \keyterm{normal equations}
$\hat{\mathbf{w}} = (\mathbf{X}^{\!\top}\mathbf{X})^{-1}\mathbf{X}^{\!\top}
\mathbf{y}$\textsuperscript{\cite{jamesIntroductionStatisticalLearning2013}} ---
subject to the numerical caution given earlier.

\subsection{Hypothesis Tests, Errors, and What a $p$-Value Is Not}

A test weighs a null hypothesis $H_0$ against an alternative. Two errors are
possible, and they trade off against each other.

\begin{figure}[htbp]
\centering
\begin{tikzpicture}
\begin{axis}[almanach axis,height=48mm,width=0.86\linewidth,
  domain=-4:7.5,samples=220,
  xlabel={test statistic},ylabel={density},
  ymin=0,ymax=0.46,ytick=\empty,xtick=\empty,
  legend style={at={(0.5,1.03)},anchor=south},legend columns=2]
  \addplot[AlmanachTeal,thick] {0.39894228*exp(-x^2/2)};
    \addlegendentry{under $H_0$}
  \addplot[AlmanachOchre,thick] {0.39894228*exp(-(x-3.4)^2/2)};
    \addlegendentry{under $H_1$}
  \addplot[fill=AlmanachRed!35,draw=none,domain=1.96:4.4]
    {0.39894228*exp(-x^2/2)} \closedcycle;
  \addplot[fill=AlmanachOchre!45,draw=none,domain=-1:1.96]
    {0.39894228*exp(-(x-3.4)^2/2)} \closedcycle;
  \draw[AlmanachInk!70,thick,dashed]
    (axis cs:1.96,0) -- (axis cs:1.96,0.42);
  \node[font=\scriptsize\sffamily,anchor=south] at (axis cs:1.96,0.42)
    {threshold};
  \node[font=\scriptsize\sffamily,color=AlmanachRed,anchor=west]
    at (axis cs:2.3,0.09) {Type I ($\alpha$)};
  \node[font=\scriptsize\sffamily,color=AlmanachOchre!75!black,anchor=east]
    at (axis cs:1.75,0.055) {Type II ($\beta$)};
\end{axis}
\end{tikzpicture}
\caption{The two error types are areas on opposite sides of one threshold.
Sliding the threshold right shrinks the red area and grows the orange one; the
only way to shrink both at once is to move the curves apart, which means a
larger sample or a bigger real effect. Statistical power is $1-\beta$.}
\label{fig:math-type-i-ii}
\end{figure}

\begin{table}[htbp]
\centering
\caption{The four outcomes of a hypothesis test. The costs of the two errors
are almost never equal, and deciding which one you can better afford is an
engineering judgement, not a statistical one.}
\label{tab:math-error-types}
\begin{tblr}{
  width=\textwidth,
  colspec={Q[l,0.20\textwidth] X[l] X[l]},
  row{1}={font=\bfseries,bg=AlmanachMist},
  hlines,vlines}
& $H_0$ actually true & $H_0$ actually false \\
\textbf{Reject $H_0$} & \textcolor{AlmanachRed}{Type I error}, rate $\alpha$
  --- a false alarm & Correct detection, probability $1-\beta$ (power) \\
\textbf{Fail to reject} & Correct non-detection, probability $1-\alpha$ &
  \textcolor{AlmanachOchre!75!black}{Type II error}, rate $\beta$ --- a
  miss \\
\end{tblr}
\end{table}

The \keyterm{$p$-value} is the probability of observing data at least as
extreme as yours \emph{assuming the null hypothesis is true}.

\begin{cautionbox}[Four things a $p$-value is not]
The American Statistical Association issued a formal statement on this because
misuse is so widespread\textsuperscript{\cite{wassersteinASAStatementPvalues2016}}.
A $p$-value is \textbf{not}:
\begin{itemize}
  \item the probability that the null hypothesis is true --- it is computed
        \emph{assuming} the null, so it cannot also measure it;
  \item the probability that your result was due to chance;
  \item a measure of effect size --- with a large enough sample, a trivially
        small difference produces $p<0.001$;
  \item a decision rule. $p = 0.049$ and $p = 0.051$ are the same evidence, and
        treating $0.05$ as a boundary between truth and falsehood has no
        statistical justification\textsuperscript{\cite{amrheinScientistsRiseUp2019}}.
\end{itemize}
Report an effect size with a confidence interval alongside any $p$-value. If
you can report only one, report the
interval\textsuperscript{\cite{reinhartStatisticsDoneWrong2015}}.
\end{cautionbox}

\begin{engineernote}
This matters directly for machine learning. ``Model B beats model A,
$p<0.05$'' says nothing about whether the improvement is large enough to
justify deployment. Ask instead: how big is the difference, how uncertain is
it, was the test set used more than once, and does the gain survive on the
slices you care about? \Cref{ch:evaluation-mlops} builds this into a protocol.
\end{engineernote}

\begin{checkpoint}
An A/B test on $2$ million users finds a conversion increase of $0.02$
percentage points with $p = 0.001$. Your colleague wants to ship it. Give two
questions you would ask before agreeing, and say which statistic you would
demand to see.
\end{checkpoint}

% ===========================================================================
\section{Deep Dive: When Correct Mathematics Gives Wrong Answers}
\label{sec:math-numerics}
% ===========================================================================

Everything above assumed real numbers. Computers use floating-point
approximations with finite precision, governed by IEEE
754\textsuperscript{\cite{ieee754-2019}}, and the gap between the two is where
a great many silent failures live\textsuperscript{\cite{goldbergWhatEveryComputer1991}}.

\subsection{Machine Epsilon and Catastrophic Cancellation}

\begin{workedexample}[Two numbers that should differ by one, and do not]
In \texttt{float32} the mantissa holds $24$ bits, so near a value $v$ the
spacing between representable numbers is about $v\times2^{-23}$.

At $v = 10^8$ that spacing is
\[
  10^8 \times 2^{-23} \approx 10^8 \times 1.19\times10^{-7} \approx 12 .
\]
So $10^8$ and $10^8+1$ round to the \emph{same} \texttt{float32} value, and
\[
  (10^8 + 1) - 10^8 = 0 \quad\text{instead of}\quad 1 .
\]
The relative error is $100\,\%$. This is \keyterm{catastrophic cancellation}:
subtracting two nearly equal large numbers annihilates every significant digit
of the answer.

\emph{Where it bites in practice.} The one-pass variance formula
$\mathbb{E}[X^2]-\mathbb{E}[X]^2$ from earlier is exactly this shape. For data
with a large mean and small variance, both terms are large and nearly equal, and
the computed variance can come out \emph{negative}. Numerical libraries use
Welford's algorithm instead --- mathematically identical, numerically stable.
\end{workedexample}

\subsection{Log-Sum-Exp}

Computing $\log\sum_i e^{x_i}$ directly overflows for moderately large inputs.
With $m = \max_i x_i$,
\begin{equation}
  \log\sum_i e^{x_i} = m + \log\sum_i e^{x_i - m}.
  \label{eq:math-logsumexp}
\end{equation}
\where{$m=\max_i x_i$ is the largest input, so every shifted exponent
$x_i-m$ is at most zero and therefore representable.}
Subtracting $m$ changes neither the exact result nor the softmax ratios, but
keeps every exponent at most zero and therefore representable.

\begin{workedexample}[The same expression, twice]
Let $x = (1000,\ 1001,\ 1002)$.

\emph{Naive.} $e^{1000}$ overflows \texttt{float64}, whose largest value is
about $1.8\times10^{308}$ and whose exponential limit is around $e^{709}$. The
computation returns \texttt{inf}, and then $\log(\texttt{inf}) =
\texttt{inf}$. Total loss.

\emph{Stabilised.} With $m = 1002$, the shifted values are $(-2,-1,0)$:
\[
  \log\bigl(e^{-2}+e^{-1}+e^{0}\bigr)
  = \log(0.1353 + 0.3679 + 1.0000)
  = \log(1.5032) = 0.4076,
\]
so the answer is $1002 + 0.4076 = 1002.4076$. Every intermediate value stayed
between $0$ and $1$.

The two expressions are algebraically identical. One returns the answer and one
returns \texttt{inf}, which is why every serious softmax implementation
subtracts the maximum first --- the point of the checkpoint at the end of
\cref{ch:deep-learning}'s softmax example.
\end{workedexample}

\subsection{Conditioning}

An algorithm can be mathematically correct and numerically unreliable. The
\keyterm{condition number} measures how much a small change in the input can
change the output.

\begin{workedexample}[A $0.05\,\%$ change in the data, a $100\,\%$ change in the answer]
Solve $\mathbf{A}\mathbf{x} = \mathbf{b}$ with
\[
  \mathbf{A} = \begin{pNiceMatrix} 1 & 1\\ 1 & 1.001\end{pNiceMatrix},
  \qquad \mathbf{b} = \begin{pNiceMatrix} 2\\ 2.001\end{pNiceMatrix}.
\]
Subtracting the first equation from the second gives $0.001x_2 = 0.001$, so
$x_2 = 1$ and $x_1 = 1$. The solution is $(1,1)$.

Now perturb the last entry of $\mathbf{b}$ from $2.001$ to $2.002$ --- a change
of $0.05\,\%$:
\[
  0.001x_2 = 0.002 \Rightarrow x_2 = 2,
  \qquad x_1 = 2 - 2 = 0 .
\]
The solution jumped from $(1,1)$ to $(0,2)$.

Nothing was computed incorrectly. The \emph{problem} amplifies input error,
because the two rows of $\mathbf{A}$ are nearly parallel --- geometrically, we
are intersecting two lines that meet at a very shallow angle, so sliding one
slightly moves the intersection a long
way\textsuperscript{\cite[][ch.~1]{highamAccuracyStabilityNumerical2002}}.
Standardising features, using orthogonal decompositions, and adding
regularisation all improve the conditioning of the problem itself rather than
merely speeding up a solver.
\end{workedexample}

\begin{cautionbox}[Every interval answers a different question]
A confidence, credible, bootstrap, prediction, or conformal interval is
computed under different assumptions and licenses a different sentence. For
every interval you report, state the random object, the conditioning
information, the coverage target, the data-dependence assumptions, and whether
any tuning reused the calibration data. Then check empirical coverage on the
slices you care about rather than trusting the nominal percentage because of
its name.
\end{cautionbox}

\begin{practicallab}[Break your own softmax]
\textbf{Goal.} Produce all three numerical failures above yourself, so you
recognise them in someone else's stack trace.

\textbf{Protocol.}
\begin{enumerate}
  \item Implement softmax naively as \texttt{exp(x)/sum(exp(x))}. Feed it
        $(1,2,3)$ and confirm it is correct. Feed it $(1000,1001,1002)$ and
        record what happens.
  \item Implement the stabilised version by subtracting $\max_i x_i$. Confirm
        both inputs now work and that the first result is unchanged to machine
        precision.
  \item In \texttt{float32}, compute $(10^8+1)-10^8$. Then compute the variance
        of $[10^8, 10^8+1, 10^8+2]$ using $\mathbb{E}[X^2]-\mathbb{E}[X]^2$ and
        compare with the two-pass formula. Note the sign of your answer.
  \item Solve the ill-conditioned system above with a library solver, perturb
        $\mathbf{b}$ in the last digit, and measure the relative change in
        $\mathbf{x}$ divided by the relative change in $\mathbf{b}$. That ratio
        is the condition number in action.
\end{enumerate}

\textbf{Deliverable.} A short table: input, naive result, stabilised result,
relative error. Keep it. When a training run produces \texttt{NaN} at step
$3000$, this table is the first place to look.
\end{practicallab}

\begin{checkpoint}
Your loss becomes \texttt{NaN} after a few hundred steps. List three
numerically distinct causes drawn from this section, and for each one name the
cheapest diagnostic that would confirm or eliminate it.
\end{checkpoint}

% ===========================================================================
\section{Integral Transforms}
% ===========================================================================

\keyterm{Convolution} combines two functions by sliding one across the other
and integrating the overlap:
\begin{equation}
  (f * g)(t) = \int_{-\infty}^{\infty} f(\tau)\,g(t-\tau)\,d\tau .
  \label{eq:math-convolution}
\end{equation}
\where{$t$ is the output position, $\tau$ the integration variable sweeping
across the overlap, and $g(t-\tau)$ the second function reversed and shifted to
$t$.}
This is the continuous parent of the discrete operation computed cell by cell
in \cref{ch:deep-learning}; the sliding filter of a convolutional network is
\cref{eq:math-convolution} with sums replacing integrals.

The \keyterm{Fourier transform} re-expresses a signal as a sum of
frequencies:
\begin{equation}
  \hat f(\omega) = \int_{-\infty}^{\infty} f(t)\,e^{-i\omega t}\,dt .
  \label{eq:math-fourier}
\end{equation}
\where{$f(t)$ is the signal in time, $\omega$ the angular frequency,
$\hat f(\omega)$ the amount of that frequency present, and $i$ the imaginary
unit.}
Its practical importance here is the \keyterm{convolution theorem}: convolution
in the time domain is multiplication in the frequency domain. That identity
underlies fast filtering, spectrogram features for speech, and the frequency
analysis used in condition monitoring.

\begin{engineernote}
When you meet a spectrogram in the speech sections of this book, remember what
it is: the signal chopped into short windows, each one passed through
\cref{eq:math-fourier}, and the magnitudes stacked as an image. That is why
audio models can borrow convolutional architectures designed for pictures ---
the transform turned a time series into a picture first.
\end{engineernote}

\subsection{Vector Calculus for Fields}

For a vector-valued $\vec{f}:\mathbb{R}^n\to\mathbb{R}^m$, the derivative is
the \keyterm{Jacobian} matrix with entries $J_{ij} = \partial f_i/\partial
x_j$\textsuperscript{\cite{bronsteinTaschenbuchMathematik2001}}:
\begin{equation}
  \mathbf{J} =
  \begin{pNiceMatrix}
    \frac{\partial f_1}{\partial x_1} & \cdots &
      \frac{\partial f_1}{\partial x_n}\\
    \vdots & \ddots & \vdots\\
    \frac{\partial f_m}{\partial x_1} & \cdots &
      \frac{\partial f_m}{\partial x_n}
  \end{pNiceMatrix}.
  \label{eq:math-jacobian}
\end{equation}
\where{$m$ is the number of outputs and $n$ the number of inputs, so
$\mathbf{J}$ is $m\times n$; row $i$ is the gradient of output $f_i$.}
The Jacobian generalises both the scalar derivative ($m=n=1$) and the gradient
($m=1$). In backpropagation, each layer contributes one Jacobian and the chain
rule multiplies them --- which is the matrix-valued form of
\cref{eq:math-chain-rule}.

A \keyterm{scalar field} assigns a number to every point (a temperature
$T(x,y,z)$); a \keyterm{vector field} assigns a vector (a wind velocity). The
gradient of a scalar field is itself a vector field pointing uphill --- exactly
the object gradient descent follows downhill.

Three further operators act on such fields\textsuperscript{\cite{bronsteinTaschenbuchMathematik2001}}:
\begin{equation}
  \operatorname{div}\vec{F} = \nabla\cdot\vec{F}
    = \sum_i \frac{\partial F_i}{\partial x_i},
  \qquad
  \operatorname{curl}\vec{F} = \nabla\times\vec{F},
  \qquad
  \Delta f = \nabla^2 f = \sum_i \frac{\partial^2 f}{\partial x_i^2}.
  \label{eq:math-div-curl-laplacian}
\end{equation}
\where{$\vec{F}$ is a vector field with components $F_i$, $f$ a scalar field,
and $\nabla$ the vector of partial derivatives; the three operators combine it
by dot product, cross product, and repetition respectively.}
Divergence measures net outward flux at a point, curl measures local rotation,
and the \keyterm{Laplacian} generalises the second derivative to several
dimensions.

\begin{engineernote}
These look like fluid dynamics, and they are --- but two of them reappear here.
The Laplacian is the operator behind diffusion, which is the physical process
that diffusion models imitate when they denoise an image step by step. Its
discrete analogue, the \keyterm{graph Laplacian} $L = D - A$ built from a
degree matrix and an adjacency matrix, is the central object in spectral graph
theory and in graph neural networks. Physics-informed networks use all three
directly, by putting the governing differential equation into the loss.
\end{engineernote}

\chapterreview{Objects}{Transform}{Estimate}{Optimise}

% ===========================================================================
\section{Further Reading}
% ===========================================================================

\begin{v3topic}{Mathematics for machine learning specifically}
\begin{itemize}
  \item \fullcite{deisenrothMathematicsMachineLearning2020} --- the closest
    match to this chapter's scope, and free online.
  \item \fullcite{strangIntroductionLinearAlgebra2016} --- the standard
    linear-algebra text; strong geometric intuition.
  \item \fullcite{boydConvexOptimization2004} --- for when you need to know
    whether an optimisation problem is genuinely hard.
\end{itemize}
\end{v3topic}

\begin{v3topic}{Probability, statistics, and information}
\begin{itemize}
  \item \fullcite{wassermanAllStatistics2004} --- concise and rigorous;
    good after this chapter rather than before it.
  \item \fullcite{mackayInformationTheoryInference2005} --- information theory
    and inference treated as one subject.
  \item \fullcite{coverElementsInformationTheory2006} --- the reference for
    entropy, KL divergence, and coding.
  \item \fullcite{brucePracticalStatisticsData2017} --- practical framing for
    data work.
  \item \fullcite{reinhartStatisticsDoneWrong2015} --- a short catalogue of the
    mistakes this chapter warns about.
\end{itemize}
\end{v3topic}

\begin{v3topic}{Numerics, and general engineering mathematics}
\begin{itemize}
  \item \fullcite{goldbergWhatEveryComputer1991} --- still the best single
    article on floating point.
  \item \fullcite{highamAccuracyStabilityNumerical2002} --- the reference for
    conditioning and stability.
  \item \fullcite{rileyMathematicalMethodsPhysics2008} --- broad coverage of
    engineering mathematics.
  \item \fullcite{bronsteinTaschenbuchMathematik2001} --- the standard
    German-language handbook; excellent as a lookup table.
  \item \fullcite{mathaiLinearAlgebraCourse2017} --- linear algebra aimed at
    physicists and engineers.
\end{itemize}
\end{v3topic}
